% ============================================================
% Preámbulo del informe final. Parte del .tex del anteproyecto
% (fuente, idioma, paquetes) y le cambia la página a ICONTEC.
% ============================================================

% ---------- Codificación e idioma ----------
\usepackage[utf8]{inputenc}
\usepackage[T1]{fontenc}
% es-tabla: "Tabla" e "Índice de tablas" en lugar de "Cuadro", que es
% como las nombran las directrices ("tablas de contenido, de figuras y
% de tablas").
\usepackage[spanish,es-nodecimaldot,es-tabla]{babel}
\usepackage{csquotes}

% ---------- Fuente tipo Calibri (la del anteproyecto) ----------
% ICONTEC solo *sugiere* Arial 12; se conserva la del anteproyecto.
\usepackage{carlito}
\renewcommand{\familydefault}{\sfdefault}

% ---------- Página según ICONTEC (NTC 1486, 6.a actualización, 2008) ----------
% Los cuatro márgenes a 3 cm, la variante simétrica que ICONTEC admite
% (la pide para impresión a doble cara), para que el texto quede
% centrado en la hoja. La variante para encuadernar es left=4cm,
% right=2cm. Número de página centrado a 2 cm del borde inferior
% (footskip = 3 cm - 2 cm); título de capítulo a 4 cm (ver titlesec).
\usepackage[letterpaper, top=3cm, left=3cm, right=3cm, bottom=3cm,
            footskip=1cm]{geometry}

% ---------- Interlineado ----------
% ICONTEC pide interlineado sencillo con una línea en blanco entre
% párrafos, que es una convención de Word. Se deja el 1.5 del
% anteproyecto, que es más legible con ecuaciones.
\usepackage{setspace}
\onehalfspacing

% ---------- Matemáticas, tablas y gráficos ----------
\usepackage{amsmath,amssymb}
\usepackage{array}
\usepackage{booktabs}
\usepackage{multirow}
\usepackage{graphicx}
\graphicspath{{figuras/}}
\usepackage{xcolor}
\usepackage{float}
\usepackage[font=small,labelfont=bf]{caption}
\usepackage{subcaption}

% ---------- Unidades ----------
\usepackage{siunitx}
\sisetup{output-decimal-marker={.}, group-digits=false}

% ---------- Listas compactas ----------
\usepackage{enumitem}
\setlist{itemsep=2pt, topsep=4pt, parsep=0pt}

% ---------- Títulos según ICONTEC ----------
% Capítulo en página nueva, centrado, en mayúscula sostenida y a 4 cm
% del borde (3 cm de margen + 1 cm).
\usepackage{titlesec}
\titleformat{\chapter}[block]
  {\normalfont\large\bfseries\centering}{\thechapter.}{0.5em}{\MakeUppercase}
\titlespacing*{\chapter}{0pt}{1cm}{1cm}
\titleformat{\section}{\normalfont\bfseries}{\thesection}{0.8em}{}
\titleformat{\subsection}{\normalfont\bfseries\itshape}{\thesubsection}{0.8em}{}
\titlespacing*{\section}{0pt}{12pt}{6pt}
\titlespacing*{\subsection}{0pt}{10pt}{4pt}

% ---------- Nombres de las listas ----------
\addto\captionsspanish{%
  \renewcommand{\contentsname}{Tabla de contenido}%
  \renewcommand{\listfigurename}{Lista de figuras}%
  \renewcommand{\listtablename}{Lista de tablas}%
}

% ---------- Bibliografía: IEEE numérico, como el anteproyecto ----------
% Para pasar a ICONTEC (NTC 5613) basta cambiar style=ieee.
% En Overleaf compila con biber sin configurar nada.
\usepackage[backend=biber, style=ieee]{biblatex}
\addbibresource{referencias.bib}

% ---------- Enlaces (va al final) ----------
\usepackage[hidelinks]{hyperref}

% ---------- Marcas de trabajo ----------
% \pendiente{...} deja en rojo lo que falta escribir. Antes de entregar,
% buscar "pendiente" en el proyecto debe dar cero resultados.
\newcommand{\pendiente}[1]{\textcolor{red}{[Pendiente: #1]}}
